\section{Proyección pública del registro interno de las rutas}

Cada ruta posee \(53\) campos en el registro unificado y \(57\) campos en el
registro enriquecido. La ficha científica publica \(49\) campos. Esta
proyección no elimina información: los estados de consistencia de la ruta,
del empalme y de la prolongación por torres, junto con las \(5\) componentes
crudas de la firma, permanecen en los registros completos y se reúnen en los
invariantes publicados. Por tanto, la aplicación de publicación es una
proyección tipada
\[
  \mathcal R_{57}\longrightarrow\mathcal F_{49},
\]
acompañada por el registro íntegro \(\mathcal R_{53}\) y por el registro
enriquecido \(\mathcal R_{57}\); no es una reducción del dominio científico.

\section{Correspondencia entre el dominio genealógico y los inventarios metrológicos}

El dominio material se construye mediante \(81\) celdas, \(56\) familias de
ruta, \(13\) multisecciones y \(324\) rutas orientadas. Su despliegue íntegro
se encuentra en el
\hyperref[chap:atlas-genealogico-exhaustivo-324-rutas]{atlas de \(324\)
rutas}, con una ficha completa y \(49\) campos por ruta. Los censos externos
de contraste se conservan en el
\hyperref[chap:inventario-universal-471-masas-384-anchuras]{inventario de
\(471\) masas y \(384\) anchuras}. El atlas construye las fibras y sus
operadores; los censos no seleccionan retrospectivamente una ruta ni
redefinen el dominio de la ley.
